% Options for packages loaded elsewhere
\PassOptionsToPackage{unicode}{hyperref}
\PassOptionsToPackage{hyphens}{url}
\documentclass[
  ukrainian,
]{article}
\usepackage{xcolor}
\usepackage[margin=2.5cm]{geometry}
\usepackage{amsmath,amssymb}
\setcounter{secnumdepth}{-\maxdimen} % remove section numbering
\usepackage{iftex}
\ifPDFTeX
  \usepackage[T1]{fontenc}
  \usepackage[utf8]{inputenc}
  \usepackage{textcomp} % provide euro and other symbols
\else % if luatex or xetex
  \usepackage{unicode-math} % this also loads fontspec
  \defaultfontfeatures{Scale=MatchLowercase}
  \defaultfontfeatures[\rmfamily]{Ligatures=TeX,Scale=1}
\fi
\usepackage{lmodern}
\ifPDFTeX\else
  % xetex/luatex font selection
\fi
% Use upquote if available, for straight quotes in verbatim environments
\IfFileExists{upquote.sty}{\usepackage{upquote}}{}
\IfFileExists{microtype.sty}{% use microtype if available
  \usepackage[]{microtype}
  \UseMicrotypeSet[protrusion]{basicmath} % disable protrusion for tt fonts
}{}
\makeatletter
\@ifundefined{KOMAClassName}{% if non-KOMA class
  \IfFileExists{parskip.sty}{%
    \usepackage{parskip}
  }{% else
    \setlength{\parindent}{0pt}
    \setlength{\parskip}{6pt plus 2pt minus 1pt}}
}{% if KOMA class
  \KOMAoptions{parskip=half}}
\makeatother
\usepackage{graphicx}
\graphicspath{{/Users/kalinicenkovitalijmikolajovic/.codex/worktrees/5c15/Script_for_translate_book/Output_Final/farlou_chunked_v1_pdf_chunks/}}
\makeatletter
\newsavebox\pandoc@box
\newcommand*\pandocbounded[1]{% scales image to fit in text height/width
  \sbox\pandoc@box{#1}%
  \Gscale@div\@tempa{\textheight}{\dimexpr\ht\pandoc@box+\dp\pandoc@box\relax}%
  \Gscale@div\@tempb{\linewidth}{\wd\pandoc@box}%
  \ifdim\@tempb\p@<\@tempa\p@\let\@tempa\@tempb\fi% select the smaller of both
  \ifdim\@tempa\p@<\p@\scalebox{\@tempa}{\usebox\pandoc@box}%
  \else\usebox{\pandoc@box}%
  \fi%
}
% Set default figure placement to htbp
\def\fps@figure{htbp}
\makeatother
\ifLuaTeX
\usepackage[bidi=basic,shorthands=off,provide=*]{babel}
\else
\usepackage[bidi=default,shorthands=off,provide=*]{babel}
\fi
\ifLuaTeX
  \usepackage{selnolig} % disable illegal ligatures
\fi
\setlength{\emergencystretch}{3em} % prevent overfull lines
\providecommand{\tightlist}{%
  \setlength{\itemsep}{0pt}\setlength{\parskip}{0pt}}
\usepackage{bookmark}
\IfFileExists{xurl.sty}{\usepackage{xurl}}{} % add URL line breaks if available
\urlstyle{same}
\hypersetup{
  pdflang={uk},
  hidelinks,
  pdfcreator={LaTeX via pandoc}}

\author{}
\date{}

\begin{document}

\pandocbounded{\includegraphics[keepaspectratio]{images/_page_316_Figure_2.jpeg}}

\textbf{РИС.} 42.3. Блок-схема розрахунку інтегрального
\(\int_{a}^{b} f(x) dx\) методом Монте-Карло (100 кидків);
\(M = \max_{a} f(x)\) .

\section{Випадкові
числа}\label{ux432ux438ux43fux430ux434ux43aux43eux432ux456-ux447ux438ux441ux43bux430}

Перш ніж перейти до рівняння Монте-Карло з частковими похідними,
зупинимося на важливому питанні випадкових чисел. Щойно, обчислюючи
інтеграл, нам потрібно було побудувати послідовність випадкових точок
\(P_i = (x_i, y_i)\) всередині прямокутника R. Іншими словами,
x-координата точки P має бути випадковим числом з інтервалу {[}a, b{]},
а \(y_i\) від інтервалу {[}0, M{]}. Щоб знайти випадкові числа з цих
інтервалів, звертаємося до послідовності випадкових чисел, рівномірно
розподілених по сегменту {[}0, 1{]}. Тоді числа, рівномірно розподілені
на сегменті {[}a, b{]}, можна обчислити за формулою

(начальное смещение струны), \(u_t\)(x, 0) = \(g(x)\) (начальная
скорость струны).

Этим волновое уравнение отличается от уравнения теплопроводности, для
которого необходимо задать только одно начальное условие.

\begin{enumerate}
\def\labelenumi{\arabic{enumi}.}
\setcounter{enumi}{3}
\tightlist
\item
  С помощью волнового уравнения можно описать распределение
  электрического тока в проводе. Из законов Кирхгофа получается
  следующая система двух уравнений с частными производными первого
  порядка:
\end{enumerate}

(16.3)

Тепер, звичайно, виникає питання: як отримати послідовність випадкових
чисел, \(\{r_i: i=1,2,\ldots\}\) рівномірно розподілених на {[}0,1{]}?
Найпоширенішим методом сьогодні є метод дедукцій (метод порівнянь).
Використовуючи цей метод, отримуємо послідовність випадкових цілих чисел
(наприклад, 2120, 1401, 177, 3013,\ldots), потім ліворуч від кожного
числа розміщується десяткова крапка, так що всі числа знаходяться між
нулем і одним \((0,212; 0,1401; 0,0177; 0,3013; \ldots)\) .

Отже, щоб отримати послідовність випадкових цілих чисел (між 0 і P), ми
використаємо метод віднімання.

\section{Отримання випадкових чисел за допомогою методу
відрахування.}\label{ux43eux442ux440ux438ux43cux430ux43dux43dux44f-ux432ux438ux43fux430ux434ux43aux43eux432ux438ux445-ux447ux438ux441ux435ux43b-ux437ux430-ux434ux43eux43fux43eux43cux43eux433ux43eux44e-ux43cux435ux442ux43eux434ux443-ux432ux456ux434ux440ux430ux445ux443ux432ux430ux43dux43dux44f.}

\begin{enumerate}
\def\labelenumi{\arabic{enumi}.}
\item
  Як перше випадкове ціле число, обирайте будь-яке число, що знаходиться
  між 0 і P (число P обирається наперед).
\item
  Помножте це випадкове число на певне фіксоване число.
\end{enumerate}

Попередньо обраний множник M.

\begin{enumerate}
\def\labelenumi{\arabic{enumi}.}
\setcounter{enumi}{2}
\item
  Додайте до добутку деяке фіксоване ціле число \emph{K}, обране
  заздалегідь.
\item
  Поділити отриману суму на P і залишок
\end{enumerate}

Виберіть нове випадкове число.

Метод виведення можна записати у формі формули
\(r_{i+1} \equiv (Mr_i + K) \bmod P\) , що означає, що для отримання
наступного \(r_{i+1}\) потрібно взяти попередній \(r_i\) , помножити
його на M, додати результат на K, поділити суму на P і взяти решту
поділу. Наприклад, якщо P = 100, M = 37, K = 16, \(r_0 = 15\) , то
отримується наступна послідовність випадкових чисел:

где

ж --- координата вдоль провода,

\emph{t} --- время,

\(i(x, t)\) --- распределение тока вдоль провода,

\begin{verbatim}
v\left(x,\ t\right)   — распределение потенциала вдоль провода,
\end{verbatim}

С -- емкость провода на единицу длины,

G --- утечка на единицу длины,

Р --- сопротивление на единицу длины,

-- индуктивность провода на единицу длины.

Уравнения (16.3) принято называть системой телеграфных уравнений. В
такой форме они обычно не используются. Для преобразования этой системы
продифференцируем первое уравнение по x, второе продифференцируем по t,
умножим обе части на C и вычтем из первого. В результате получаем

\[
Mr_1 + K = 2643
\]

Воспользуемся теперь вторым уравнением из (16.3)

\[
r = 71
\]

окончательно получаем

(16.4)

Y

Это уравнение для тока, известное под названием \emph{телеграфного
уравнения}, является гиперболическим уравнением второго порядка (если
конечно ии C, ни L не равны нулю, в противном случае уравнение
становится параболическим).

Напряжение описывается точно таким же уравнением:

(16.5)

\[
I = \int_{0}^{1} \int_{0}^{1} \int_{0}^{1} \int_{0}^{1} e^{-(x^{2}+y^{2}+z^{2}+\omega^{2})} dx dy dz dw,
\]

Если G = R = 0, то уравнения (16.4)---(16.5) упрощаются:

(16.6)

\[
I = \int_{0}^{1} e^{\sin x} dx.
\]

\section{ЗАДАЧИ}\label{ux437ux430ux434ux430ux447ux438}

\begin{itemize}
\tightlist
\item
  \begin{enumerate}
  \def\labelenumi{\arabic{enumi}.}
  \tightlist
  \item
    Получите уравнение (16.5) для \emph{и} из системы двух уравнений
    первого порядка (16.3)
  \end{enumerate}
\item
  \begin{enumerate}
  \def\labelenumi{\arabic{enumi}.}
  \setcounter{enumi}{1}
  \tightlist
  \item
    Зная физический смысл каждого члена волнового уравнения, что вы
    можете сказать о поведении во времени решения следующей задачи:
  \end{enumerate}
\end{itemize}

(YYII)

\[
r_{i+1} \equiv (3r_i + 4) \mod 7, \quad r_0 = 0.
,   0 < x < 1  ,   0 < t < \infty  ,
\]

(ГУ)

\[
I = \int_{0}^{1} \int_{0}^{1} \int_{0}^{1} e^{-(x^{2}+y^{2}+z^{2})} dx dy dz.
u(x,y) = g(x,y) =
\begin{cases} u(0, t) = 0, & 0 < t < \infty, \\ u(1, t) = 0, & 0 < t < \infty, \end{cases}
\]

\$\$ После того как эта частная задача будет решена методом Монте-Карло,
мы обсудим, как решать этим методом более общие задачи.

Для иллюстрации метода Монте-Карло рассмотрим игру, которая называется
«Блуждающий пьяница». Чтобы играть в эту игру, необходим стол с
нанесенной на нем сеткой из линий, как показано на рис. 43.1.

Перейдем к описанию правил игры.

\end{document}
